\subsection*{一步一步建立直观}

\paragraph{1. 用同一个初值锁定两种表达}

余弦和正弦组成的位置 $(\cos t,\sin t)$ 从 $(1,0)$ 出发，速度是位置旋转直角后的向量。复指数 $e^{it}$ 满足同样的局部规则与初值。若调用相应微分方程的唯一性，两种表达就描述同一运动，而不只是偶然在几个角度上数值相同。



\paragraph{2. 半圈的位置给出恒等式}

当时间走到半圈，横坐标为负一、纵坐标为零，因此 $e^{i\pi}=-1$。右图保留速度箭头，强调等式中的常数来自整条连续运动在特殊时刻的位置。它不是把五个符号强行拼在一起，而是指数结构与圆周几何的交汇。



\paragraph{3. 检查哪一种修改会破坏单位圆}

把实部增长率调正，点仍旋转，但半径不再固定。此时半圈位置不再是负一，而是负实轴上一个更远的点。这个对照帮助检验欧拉公式的准确条件：纯虚指数只改变相位，加入实部还会乘上实指数的模长因子。


